% paper/sections/conclusion.tex

\section{Conclusion}
\label{sec:conclusion}

We presented the first systematic study of curriculum domain
randomisation for \AUV{} fluid physics, comparing \CDR{} against
Uniform \DR{}, Naive \SAC{}, and a \PID{} baseline across nine
experiments.
Three findings have practical significance for \AUV{} deployment.

First, \PID{} control collapses catastrophically
($100\% \to 3\%$ success) under modest fluid-parameter shift,
rendering fixed-gain classical control unsuitable for real ocean
deployment without real-time parameter adaptation.

Second, domain randomisation is essential for reliable transfer:
it eliminates the $\pm 33.83$\,pp inter-seed variance
of unrobust policies, enabling confident pre-deployment performance
assessment.

Third, \CDR{} produces $1.13\%$ more energy-efficient policies
than Uniform \DR{}, a meaningful operational advantage for
battery-constrained \AUV{} missions where thrust consumption limits
mission duration.

Beyond these results, this work establishes energy efficiency as a
meaningful differentiator between \DR{} strategies that achieve
similar success rates.
As \AUV{} \RL{} matures toward physical deployment, thrust economy
should join success rate as a standard evaluation metric.
We open-source the Halcyon~X4 environment and training pipeline to
support this.